\section{Weak residual identities and discrete differentiation}
\label{sec:residuals}

We derive the residual identities needed for the two posterior estimates.
Moments of a fixed finite sequence of primary iterates will locate two
consecutive small increments. Ordinary test-vector recipes then control
means under square reweighting; one additional implicit recipe supplies
the linear-response test for the directional covariance estimate. The
proof proceeds from discrete derivative bounds to a weak residual rule,
and ends with the terminal identity for that implicit test.

This section is deterministic.  We fix $J$ satisfying the word estimates of
Lemma~\ref{lem:words}, fix an external field $h_0\in\R^n$, and take all
expectations with respect to $\mu_{h_0}$. We fix the depth, the finite
construction, and
its coefficient bounds independently of $n$. These choices determine a
finite required word length, which can increase when the choices change.
Constants
in derivative estimates can depend on these choices and on the corresponding
word bounds; they are uniform in $n$, $h_0$, the spin configuration, and the
permitted seeds.  The size and inverse bounds for the initial implicit
construction below have the stronger, depth-independent uniformity explicitly
stated there.

We use the discrete derivatives, product rule
\eqref{eq:discrete-product}, and norm $\norm G_*$ introduced above,
with $dg=(d_i g)_{i=1}^n$ for scalar $g$.  All vector functions such as
$\tanh h^1$ are understood coordinatewise.

\subsection{Primary fields}

The following recursion has the Onsager-corrected form of iterative TAP
constructions, as in Bolthausen~\cite{Bolthausen2014}. Here it starts from
a Gibbs spin sample and is used only through a fixed finite depth.
Define
\begin{equation}
 \begin{gathered}
 m^0=x,\quad b_0=0,\qquad
 m^l=\tanh h^l,\quad b_l=\av{1-(m^l)^2},\quad R_l=m^{l-1}-m^l,\\
 h^{l+1}=h_0+Jm^l-jb_lm^{l-1}\quad(l\geq1).
 \end{gathered}
 \label{eq:amp}
\end{equation}
The finite-depth objective is visible in the scalar pairings
$S_l=R_l^{\mathsf T}m^l/\sqrt n$. Expanding the squared increment gives
\begin{equation}\label{eq:residual-telescoping}
 \frac{\norm{R_l}^2}{n}=b_l-b_{l-1}-\frac{2S_l}{\sqrt n}.
\end{equation}
Since $0\le b_l\le1$, control of $S_l/\sqrt n$ through a sufficiently
large fixed depth forces two consecutive residuals to be small. We will
prove the moment bounds
\[
 \E S_l^2\le C_l,\qquad
 \E_\nu\left(\frac{S_l}{\sqrt n}\right)^2
 \le\frac{C_l}{\sqrt n}\left(1+\frac{\D_{h_0}(f)}{N}\right),
 \quad N=\E f^2>0,\quad \nu=\frac{f^2\mu_{h_0}}{N}.
\]
These bounds will supply the ordinary and square-reweighted selection
estimates in Section~\ref{sec:consequences}; convergence of the primary
recursion is not required.

The first residual explains what is needed from a test vector $U$.
Conditional covariance in one spin gives the exact identity
\[
 \E G R_1^{\mathsf T}U=\E\sum_i v_i^1d_i(GU_i).
\]
The product rule bounds its right side by $C\norm G_*$ whenever
$\norm U$ and $\sum_i|d_iU_i|$ are bounded pointwise. Thus the essential
property of our tests is a bound on the sum of their diagonal derivatives,
in addition to a bound on their Euclidean norm. The Onsager corrections
will preserve this property through the finite constructions below.
We begin with the primary fields, whose derivatives enter those tests.
We use
$\norm{Q}_{\ell^4}=(\sum_{i,j}|Q_{ij}|^4)^{1/4}$ for the unnormalized
entrywise norm.

\begin{lemma}[Primary differentiation]\label{lem:primary-derivatives}
Let
\begin{equation}
 H_1=J,\qquad M_0=I,\qquad
 M_l=D_{1-(m^l)^2}H_l,\qquad
 H_{l+1}=JM_l-jb_lM_{l-1}.
 \label{eq:primary-formal}
\end{equation}
For each fixed $l\geq1$,
\begin{align}
 &\norm{\nabla h^l-H_l}_{\Frob}
   +\norm{\nabla m^l-M_l}_{\Frob}\leq C,\label{eq:primary-error}\\
 &\norm{\nabla h^l}+\norm{\nabla m^l}
   +\norm{\nabla h^l}_{\ell^4}
   +\norm{\diag\nabla h^l}_2\leq C,
   \qquad \norm{db_l}\leq Cn^{-1/2}.
 \label{eq:primary-bounds}
\end{align}
More generally, let $g_i$ be a smooth function of a fixed list of coordinates
$h_i^l$ and averages $b_l$, with uniformly bounded derivatives through the
orders being used.  Deterministic site labels are permitted, provided the
same bounds hold uniformly in $i$.  Its derivative matrix has bounded
operator norm, bounded diagonal $\ell^2$ norm, and bounded entrywise
$\ell^4$ norm.
\end{lemma}

\begin{proof}
We first bound the formal matrices in \eqref{eq:primary-formal}, then
compare them with the actual discrete derivatives. Expanding $H_{l+1}$
produces the chain
\[
 J D_{1-(m^l)^2}J\cdots D_{1-(m^1)^2}J
\]
and the terms obtained by contracting disjoint pairs of originally adjacent
$J$'s.  Contracting $J D_{1-(m^s)^2}J$ contributes
$-j b_s I$.  This description follows inductively from the two terms defining
$H_{l+1}$: either keep its first link, or contract that link with the next
one.  Multiplicities are retained when different contractions happen to give
the same matrix.

The contraction terms cancel the diagonal predictions. To see this, fix
a complete noncrossing pairing of the original chain.  Reinserting each
contracted adjacent pair identifies its contribution with the prediction of
this full pairing, multiplied by the signs of the contractions.  The
originally adjacent pairs belonging to a pairing are mutually disjoint, and
every nonempty complete noncrossing pairing contains at least one such pair.
Summing over their subsets gives $(1-1)^a=0$, where $a\geq1$ is their
number.  An odd chain has no complete pairing.  Lemma~\ref{lem:words}
therefore bounds the diagonal $\ell^2$ norm of $H_l$, and bounds its operator
and off-diagonal $\ell^4$ norms term by term.  Its full $\ell^4$ norm is
bounded as well, since $\norm{\diag H_l}_4\leq\norm{\diag H_l}_2$.

We now compare these formal bounds with the actual derivatives. For a
scalar function $\psi$ with bounded second derivative, Taylor's theorem gives
\[
 \abs{d_j\psi(h_i^l)-\psi'(h_i^l)d_jh_i^l}
 \leq C|d_jh_i^l|^2.
\]
The Frobenius norm of the remainder is at most
$C\norm{\nabla h^l}_{\ell^4}^2$.  This proves the magnetization part of
\eqref{eq:primary-error} once the field part is known.  The same argument
for $\psi(z)=\sech^2z$ gives an operator bound for its site derivative
matrix $Q_l$.  Averaging its rows yields
\[
 db_l=n^{-1}Q_l^{\mathsf T}\one,\qquad
 \norm{db_l}\leq n^{-1/2}\norm{Q_l}.
\]
In differentiating $b_lm^{l-1}$, the term from $db_l$ is a rank-one
matrix of Frobenius norm at most
$\norm{m^{l-1}}\norm{db_l}\leq C$.  The discrete product error is
$\nabla m^{l-1}$ times a diagonal whose diagonal has norm at most
$2\norm{db_l}$.  Its Frobenius norm is bounded, since
$\norm{\nabla m^{l-1}}\leq C$.  Starting at
$\nabla h^1=J$, these estimates close the induction. The resulting bounded
Frobenius perturbation preserves the operator, diagonal $\ell^2$, and
entrywise $\ell^4$ bounds of the formal matrices.

For the final assertion, apply the same comparison to a smooth function
of several arguments. Expand along the segment between the arguments at
$x$ and $x^{(j)}$. The preceding bounds control the primary linear terms.
A linear term involving $db_s$ is a
rank-one matrix with a bounded-entry left vector, hence has bounded
Frobenius norm.  The squared-average remainder is controlled by
\[
 n\sum_j |d_jb_s|^4\leq n\norm{db_s}^4\leq C.
\]
Mixed remainders are bounded by sums of pure squares, and primary quadratic
remainders are controlled by the entrywise $\ell^4$ bounds.  This proves
the assertion.
\end{proof}

We have in particular bounded the Euclidean change of every primary
vector under a single spin flip, at each fixed depth. This is the estimate
that will keep the local constructions inside their buffers.

\subsection{Finite recipes for small vectors}

The residuals will be tested against vectors built by finitely many
operations. We call these vectors ``small'' when their Euclidean norm is
bounded; the word does not refer to coordinatewise magnitude. The following
definition specifies the operations and the meaning of their site partials.

\begin{definition}[Ordinary admissible recipe]\label{def:small-recipe}
Choose finitely many deterministic seeds $s\in\R^n$ of bounded Euclidean
norm.  Seeds may depend on $J$ and $h_0$, but are fixed as functions of $x$.
Choose a finite tuple of scalar functions $\theta$ satisfying
\[
 \sup_x\bigl(|\theta_\alpha(x)|+\norm{d\theta_\alpha(x)}\bigr)\leq C
\]
for each component.  Such scalar functions will be called parameters; they
may include the averages $b_l$.

Given auxiliary fields $y^1,\ldots,y^{a-1}$ already constructed, a source
has the form
\begin{equation}
 w_i^a=\sum_{s}\phi_{a,s,i}\bigl((h_i^l)_{l\in I_a},\theta\bigr)s_i
       +\sum_{b<a}\phi_{a,b,i}\bigl((h_i^l)_{l\in I_a},\theta\bigr)y_i^b.
 \label{eq:source-form}
\end{equation}
All sums and index sets are finite.  Coefficients are smooth and bounded,
with bounded derivatives of every fixed order needed in the construction,
uniformly on their argument ranges and the connecting line segments.
Define the new auxiliary field by
\begin{equation}
 y^a=Jw^a
       -j\sum_l\av{\partial_lw^a}\,m^{l-1}
       -j\sum_{b<a}\av{\partial_{y^b}w^a}\,w^b.
 \label{eq:small-field}
\end{equation}
The partial $\partial_l$ differentiates the explicitly displayed local
argument $h_i^l$, holding every parameter and every auxiliary argument fixed.
In particular it holds $b_l$ fixed even when that average is among the
parameters.  The partial $\partial_{y^b}$ differentiates the indicated
linear auxiliary argument, with its coefficient held fixed.  A partial in
an absent argument is zero.  Thus $\partial_{y^b}w^a$ is simply the
coefficient vector $\phi_{a,b}$.

A terminal vector is another source of the form \eqref{eq:source-form},
using any already constructed auxiliary fields.  A source $w^b$ used again
as a terminal vector or in a new source is replaced by its expression
\eqref{eq:source-form}; auxiliary fields themselves are not expanded.
The recipe is \emph{admissible from level $p$} if every explicit primary
argument appearing in its source formulas or terminal formula has index
at least $p$.  This condition does not restrict the parameters, or the
vectors $m^{l-1}$ in the correction terms of \eqref{eq:small-field}.
\end{definition}

The definition permits two operations that we will use repeatedly.
Multiplication by a bounded smooth coefficient of the same allowed
arguments gives another source. At any finite stage we may also add
parameters satisfying the stated size and difference bounds.  We will use
the deterministic seed $\one/\sqrt n$ freely; thus $m^l/\sqrt n$ is a
source with the single primary argument $h^l$.  Linearity in the
auxiliary arguments is preserved in every source. It will be used in both
the derivative estimates and the trace cancellation.

For a terminal vector $U$ admissible from level $p$, the estimate we
seek is
\[
 \left|\E G R_p^{\mathsf T}U\right|\le C\norm G_*.
\]
Its square-density counterpart bounds
$|\E f^2H_0R_p^{\mathsf T}U|$ by
$C(\E f^2+\D_{h_0}(f))$ when $H_0$ has bounded size and difference norm.
With $U=m^p/\sqrt n$, the choices $G=S_p$ and $H_0=S_p/\sqrt n$
will give the moment estimates above. Lemma~\ref{lem:residual} proves
these rules after we establish the required diagonal derivative bounds.

\subsection{A scalar identity for comparing fields}

To pass from $R_{p+1}=\tanh h^p-\tanh h^{p+1}$ to earlier residuals,
we use a scalar identity valid for arbitrary real fields. For
$z,r_0,a\in\R$, define
\begin{equation}
 \Phi(z;r_0,a)=\int_0^1 e^{a(1-u^2)/2}
       \frac{\cosh(u(z-r_0))}{\cosh z\cosh r_0}\,\dd u.
 \label{eq:phi}
\end{equation}
Differentiating $e^{a(1-u^2)/2}\sinh(u(z-r_0))$ in $u$ gives
\begin{equation}
 \bigl[(z-r_0-a\tanh z)-a\partial_z\bigr]\Phi(z;r_0,a)
       =\tanh z-\tanh r_0.
 \label{eq:phi-stein}
\end{equation}
The boundary at $u=0$ is zero; the boundary at $u=1$, after division
by $\cosh z\cosh r_0$, is the right side. The identity holds for either
sign of $a$. At $a=0$, $\Phi$ is the secant slope of $\tanh$, with
continuous value $\sech^2r_0$ at $z=r_0$.

These formulas also give the uniform coefficient bounds needed for the
recipes. Positivity of the integrand and the secant-slope formula imply
$0<\Phi\le e^{|a|/2}$. For $a$ in a fixed bounded interval, every
fixed-order derivative in $z,r_0,a$ is bounded by a constant times $\Phi$:
derivatives of the hyperbolic factors introduce bounded ratios, and those
of the exponential introduce bounded powers of $(1-u^2)/2$.
The quotient rule then bounds every fixed-order derivative of each ratio
$(\partial^\alpha\Phi)/\Phi$. These assertions are uniform in the two
real field arguments and require no stability assumption.

In the ordinary residual induction we will substitute
$(z,r_0,a)=(h_i^p,h_i^{p+1},j(b_p-b_{p-1}))$. The same identity also
provides the two tests comparing an approximate field with the stable
root, which we construct next.

\subsection{An initial implicit auxiliary field}

The directional estimate needs one auxiliary field defined by a linear
implicit equation. We construct it near the stable root, establish its
inverse margin, and then control its discrete derivatives on a smaller
region.

Assume the stable-field implication \eqref{eq:stable-field} holds at $h_0$
through the fixed residual threshold required in Lemma~\ref{lem:root}.
Let $r$ be its unique root, and put
$t_*=\tanh r$ and $q_*=\av{t_*^2}$.  Fix $k\geq2$ and consider
\begin{equation}
 \Omega_\rho=
 \{x:\max(\norm{R_k(x)},\norm{R_{k-1}(x)})\leq\rho\sqrt n\}.
 \label{eq:small-region}
\end{equation}
All choices of $\rho$ below depend only on the stable-field constants and
the operator bound on $J$, until derivative estimates are considered.
For sufficiently small $\rho$, the primary recursion and
Lemma~\ref{lem:root} give on $\Omega_{2\rho}$
\begin{equation}
 \norm{h^k-r}+\norm{m^{k-1}-t_*}\leq C\rho\sqrt n,
 \qquad |1-b_{k-1}-q_*|\leq C\rho.
 \label{eq:implicit-near-root}
\end{equation}
The comparison is independent of $k$: the primary recursion gives
\[
 F(h^k)=JR_k-jb_{k-1}(R_{k-1}+R_k)
                     +j(b_k-b_{k-1})m^k,
 \qquad
 |b_k-b_{k-1}|\leq2\norm{R_k}/\sqrt n,
\]
so the residual is at most $C\rho\sqrt n$. Applying the root estimate
therefore gives the asserted comparison with the same depth independence.

Use the scalar kernel with
\begin{equation}
 a=j(1-b_{k-1}-q_*),\qquad A_i=\Phi(h_i^k;r_i,a).
 \label{eq:phi-data}
\end{equation}
Write $A=D_{(A_i)}$ and $\chi=\av{A_i}$.  The initial source $w_0$
is one of
\[
 w_{0,i}=A_i e_i\quad(\norm e\leq1),\qquad
 w_{0,i}=n^{-1/2}\partial_{r_0}\Phi(h_i^k;r_i,a).
\]
The partial in the second formula holds $a$ fixed. The two choices
have complementary roles. Holding $r_i,a$ fixed in the site partial
$\partial_k$, the scalar identity and its derivative in $r_0$ give
\begin{align}
 \sum_i[(h_i^k-r_i-am_i^k)-a\partial_k](A_i e_i)
   &=(m^k-t_*)^{\mathsf T}e,\notag\\
 \sum_i[(h_i^k-r_i-am_i^k)-a\partial_k]
   \frac{\partial_{r_0}\Phi(h_i^k;r_i,a)}{\sqrt n}
   &=\sqrt n\{\chi-(1-q_*)\}.
 \label{eq:implicit-source-actions}
\end{align}
Indeed differentiating \eqref{eq:phi-stein} in $r_0$ gives
$[(z-r_0-a\tanh z)-a\partial_z]\partial_{r_0}\Phi
=\Phi-\sech^2r_0$. Thus the first source tests a direction of the
magnetization error, while the second tests the discrepancy of the
averaged coefficient $\chi$. Both quantities enter the posterior
comparison. To put these scalar identities into a form controlled by the
weak residual rule, introduce the coupled source and field
\begin{equation}
 \begin{gathered}
 w=w_0+A(y^{\rm imp}-jc_1t_*),\qquad
 y^{\rm imp}=Jw-j\chi w-jc_1m^{k-1},\\
 c_1=\av{\partial_k w}.
 \end{gathered}
 \label{eq:implicit}
\end{equation}
In this partial, $y^{\rm imp}$, $c_1$, and $a$ are independent arguments
held fixed.  Consequently $\partial_{y^{\rm imp}}w=A$, and the field
equation in \eqref{eq:implicit} is exactly \eqref{eq:small-field} with
one allowed self-dependence: $\chi$ is its auxiliary-partial average,
and $c_1$ its primary-partial average. The correction $jc_1m^{k-1}$
will supply the scalar cancellation in Lemma~\ref{lem:terminal-residual}.
Its only primary local argument is $h^k$.
Once constructed, $y^{\rm imp}$ is available as the first auxiliary field
in any ordinary finite recipe.  The scalar $\sqrt n c_1$ is then an
additional parameter, and $t_*/\sqrt n$ is an additional seed.

\begin{lemma}[Implicit construction]\label{lem:implicit-construction}
For sufficiently small $\rho>0$, \eqref{eq:implicit} has a unique solution
on $\Omega_{2\rho}$.  There are constants independent of the primary
depth $k$ such that, throughout this region,
\begin{equation}
 \norm w+\norm{y^{\rm imp}}+\sqrt n|c_1|\leq C,
 \qquad
 I+j\chi A-A^{1/2}JA^{1/2}\succeq cI.
 \label{eq:implicit-size}
\end{equation}
On $\Omega_{3\rho/2}$, for all sufficiently large $n$ depending on the
fixed depth, the solution also satisfies
\begin{equation}
 \norm{\nabla w}_{\Frob}+\norm{\nabla y^{\rm imp}}_{\Frob}
       +\sqrt n\norm{dc_1}\leq C_k.
 \label{eq:implicit-frobenius}
\end{equation}
\end{lemma}

\begin{proof}
We first control the coefficients and solve the implicit linear system.
Only then do we differentiate it, using the larger region to cover spin
flips from the smaller one. The uniform kernel and quotient bounds proved
above will control the coefficients even when some $A_i$ are small.

Put $\ell_i=(\partial_z\Phi/\Phi)(h_i^k;r_i,a)$.
At $z=r_0$ this ratio equals $-\tanh r_0$, for every $a$.  Comparison
with $z=r_0,a=0$, followed by \eqref{eq:implicit-near-root}, gives
\begin{equation}
 \norm{A-D_{1-(m^k)^2}}_{\Frob}\leq C\rho\sqrt n,
 \qquad \norm{\ell+t_*}\leq C\rho\sqrt n,
 \qquad |\chi-b_k|\leq C\rho.
 \label{eq:implicit-coefficients}
\end{equation}
In particular, the comparisons put $A$ among the diagonals allowed by
the stable-field implication, once $\rho$ is decreased.

Let $u_0=w_0/A$ coordinatewise and $p_0=A\partial_k u_0$.  The quotient
bounds just proved give $\norm{u_0}+\norm{p_0}\leq C$; their local
derivatives have coordinatewise bounds proportional to $|e_i|$ in the
first case and to $n^{-1/2}$ in the second.  Direct differentiation with
the specified partial convention gives
\[
 \partial_k w=D_\ell w+p_0,\qquad
 c_1=\av{\ell w+p_0}.
\]
We can now eliminate $y^{\rm imp}$ and $c_1$ and solve for the source
alone. The resulting linear equation is
\begin{equation}
 \left[I-A(J-j\chi I)
       +\frac jn A(m^{k-1}+t_*)\ell^{\mathsf T}\right]w
 =A\left[u_0-j(m^{k-1}+t_*)\av{p_0}\right].
 \label{eq:implicit-linear-system}
\end{equation}
Write the matrix on the left as $I+AL_2$.  Its symmetrically normalized
version $I+A^{1/2}L_2A^{1/2}$ differs in operator norm by $O(\rho)$ from
\[
 I+jb_k A-A^{1/2}JA^{1/2}
       -\frac{2j}{n}A^{1/2}m^k(m^k)^{\mathsf T}A^{1/2}.
\]
For the rank term, use \eqref{eq:implicit-near-root} and
\eqref{eq:implicit-coefficients} in the elementary inequality
$\norm{uv^{\mathsf T}-u'v'^{\mathsf T}}
\leq\norm{u-u'}\norm v+\norm{u'}\norm{v-v'}$.
The displayed symmetric matrix has the fixed positive margin supplied
by \eqref{eq:stable-field}. Decrease $\rho$ until the perturbation is less
than half that margin. The normalized matrix then has bounded inverse,
even though it need not itself be symmetric.  The identity
\[
 (I+AL_2)^{-1}
 =I-A^{1/2}(I+A^{1/2}L_2A^{1/2})^{-1}A^{1/2}L_2
\]
gives the same bound for the unnormalized inverse, since $\norm{L_2}$ is
bounded.  The right side of \eqref{eq:implicit-linear-system} is bounded
in norm, because $|\av{p_0}|\leq\norm{p_0}/\sqrt n$.
This establishes the size bounds. To obtain the last assertion of
\eqref{eq:implicit-size}, remove the negative rank term from the stable-field
matrix and change $b_k$ to $\chi$. The constants used to solve the system
and obtain these bounds depend only on the fixed stability and operator
margins.

It remains to bound the discrete derivatives. Compare the equations at
$x$ and $x^{(i)}$. On the stated inner region both configurations belong
to $\Omega_{2\rho}$ for large $n$, by
Lemma~\ref{lem:primary-derivatives}. In every coefficient-difference term,
retain the solution vector at the common point $x$.  In particular,
\begin{align*}
 d_i c_1
 &=\av{(d_i\ell)w(x)+d_i p_0}
             +n^{-1}\ell(x^{(i)})^{\mathsf T}d_iw,\\
 d_i\{c_1(m^{k-1}+t_*)\}
 &=c_1(x)d_i m^{k-1}
               +(m^{k-1}(x^{(i)})+t_*)d_ic_1.
\end{align*}
Use the same product arrangement for $A$ and $\chi$ in
\[
 w=A\{u_0+(J-j\chi I)w-jc_1(m^{k-1}+t_*)\}.
\]
For clarity, put $v=m^{k-1}+t_*$ and
$\mathcal B=u_0+(J-j\chi I)w-jc_1v$, and use primes for evaluation at
$x^{(i)}$ in the next display.  With
$\zeta_i=n^{-1}\{(d_i\ell)^{\mathsf T}w+\one^{\mathsf T}d_ip_0\}$,
the exact equation is
\begin{align}
 &\left[I-A'(J-j\chi'I)+\frac jn A'v'(\ell')^{\mathsf T}\right]d_iw
 \notag\\
 &\qquad=(d_iA)\mathcal B+A'd_iu_0
       -jA'w\,d_i\chi-jc_1A'd_im^{k-1}-jA'v'\zeta_i.
 \label{eq:implicit-flip-system}
\end{align}
Thus the left matrix is precisely the matrix in
\eqref{eq:implicit-linear-system} evaluated at $x^{(i)}$. The bracket
$\mathcal B$ on the right is consequently one common-point
vector for all columns. This allows the following estimates to control
the Frobenius norm of the entire right side.

We bound its terms in turn.  A site coefficient difference
times a common vector $s$ of bounded norm has bounded Frobenius norm,
since
\[
 \sum_{i,j}s_i^2|d_j\psi_i|^2
 \leq \norm s^2\max_i\sum_j|d_j\psi_i|^2\leq C.
\]
This applies to $A,\ell,u_0,p_0$, using their coefficient bounds and
Lemma~\ref{lem:primary-derivatives}; the seed factors are retained for
$u_0,p_0$.  The bracket multiplying $d_iA$ is bounded in norm by the
already established size estimates, so no bound on $1/A_i$ is needed.
Also $\norm{d\chi}\leq C/\sqrt n$, because $A_i$ is a primary site
function with an average parameter of that derivative scale.  If
$\zeta_i=\av{(d_i\ell)w+d_i p_0}$, then
\[
 \norm\zeta\leq n^{-1/2}
   \bigl(\norm{D_w\nabla\ell}_{\Frob}
                      +\norm{\nabla p_0}_{\Frob}\bigr)
 \leq C/\sqrt n.
\]
Its multipliers $m^{k-1}(x^{(i)})+t_*$ have norm at most $2\sqrt n$,
which still gives a bounded Frobenius norm after summing the squared
column norms.  Finally
$|c_1|\norm{\nabla m^{k-1}}_{\Frob}\leq C$, by its operator bound and
$|c_1|\leq C/\sqrt n$.  Applying the uniformly bounded flipped inverses
column by column proves $\norm{\nabla w}_{\Frob}\leq C_k$.
Returning to the two displayed difference identities gives
$\norm{dc_1}\leq C_k/\sqrt n$ and the required bound for
$\nabla y^{\rm imp}$.
The operator in \eqref{eq:implicit-linear-system}, the regions
$\Omega_\rho$, and the required buffer inclusions do not depend on the
directional seed $e$.  All bounds for its right side and its derivatives
use only $\norm e\leq1$; indeed $u_0=e$ and $p_0=0$ for that source.
The inverse margins, the threshold on $n$, and all these estimates are
therefore uniform over the full Euclidean unit ball of seeds.
\end{proof}

\paragraph{Local recipes and their buffers.}
A local recipe starts with this implicit field and adds only ordinary
auxiliaries. Such
recipes are evaluated on the inner
region $\Omega_\rho$, with all estimates also available on its fixed
finite flip-neighborhoods.  To define them as functions on the whole cube,
keep the actual solution through $\Omega_{2\rho}$ and assign finite values
to $y^{\rm imp},c_1$ outside that region, defining sources by their stated
expressions. The implicit equation is required only inside the buffer.
Every test
in a local statement is supported on $\Omega_\rho$. For a fixed recipe,
once $n$ is sufficiently large, every difference used in the proof stays
strictly inside the buffer. The arbitrary extension therefore does not
enter the estimates.

The buffer sizes are chosen before the construction depth; the threshold
on dimension is chosen afterward. To verify this order, choose a finite
recipe and let $q$ be the largest number of successive spin flips used by
any difference in its proof. The induction below adds only finitely many
operations at each smaller residual index, so $q$ is finite. The primary
bounds give, for every involved level $l$,
\[
 \|R_l(x^{(i)})-R_l(x)\|
 \le 2\|\nabla R_l(x)e_i\|\le C_l.
\]
With $C=\max_l C_l$, a path of at most $q$ flips increases either
residual defining $\Omega_\rho$ by at most $qC$. Thus
$n>(2qC/\rho)^2$ puts every required neighborhood of
$\Omega_{3\rho/2}$ strictly inside $\Omega_{2\rho}$. Here $q,C$, and the
threshold on $n$ may depend on the already fixed recipe and depth;
$\rho$ still depends only on the earlier stability and size margins.

\subsection{Trace control for small derivatives}

The primary derivative bounds do not yet give the diagonal trace control
needed for integration by parts. For small vectors, the source structure
and its correction terms supply this additional cancellation. We prove
it together with the size and derivative bounds used to construct later
auxiliary fields.

\begin{lemma}[Small-vector differentiation]\label{lem:small-derivatives}
Every source, auxiliary field, and terminal vector $U$ in a fixed ordinary
admissible recipe satisfies
\begin{equation}
 \norm U+\norm{\nabla U}_{\Frob}
                   +\norm{\diag\nabla U}_1\leq C.
 \label{eq:small-derivative-bound}
\end{equation}
The same holds locally for recipes with the single initial implicit field
above.  In these recipes the initial partial $\partial_k w$ is also a
small source.  For every primary coefficient and auxiliary coefficient
in \eqref{eq:small-field}, respectively,
\begin{equation}
 \begin{split}
 \alpha_l=\av{\partial_lw^a}:&\quad
       |\alpha_l|+\norm{d\alpha_l}\leq C/\sqrt n,\\
 \gamma_b=\av{\partial_{y^b}w^a}:&\quad
       |\gamma_b|+\norm{d\gamma_b}\leq C.
 \end{split}
 \label{eq:small-average-bounds}
\end{equation}
All local statements include the required fixed flip buffers.  Constants
can depend on the entire finite recipe and its primary depth.
\end{lemma}

\begin{proof}
The proof has two parts. We bound the errors introduced by discrete
differentiation, then cancel the trace predictions of the retained formal
terms. Write $\norm E_{\rm nuc}$ for the sum of the singular values of
$E$. The error estimates repeatedly use
\begin{equation}
 |\Tr(PE)|\leq\norm P\norm E_{\rm nuc},\qquad
 \norm{BC}_{\rm nuc}\leq\norm B_{\Frob}\norm C_{\Frob},\qquad
 \norm E_{\rm nuc}\leq\sum_j\norm{E_{\cdot j}}.
 \label{eq:nuclear-tools}
\end{equation}
Every bounded-length diagnostic word $P$ has bounded operator norm and
bounded off-diagonal entrywise $\ell^4$ norm. For the errors arising before
we solve the differentiated implicit
equations, $P$ may contain one inverse $K$ from \eqref{eq:K}. Its required
margin is supplied by the last assertion of \eqref{eq:implicit-size}.

\paragraph{Source errors.}
Consider one source term $\phi_i s_i$.  Retain the derivative
\begin{equation}
 D_\phi\nabla s+\sum_lD_{s\partial_l\phi}\nabla h^l,
 \label{eq:retained-source}
\end{equation}
and put parameter derivatives into the error rather than into the formal
expression. Set
\[
 Q_{ij}=\sum_l|d_jh_i^l|,\qquad
 \Theta_j=\sum_\alpha|d_j\theta_\alpha|.
\]
The finite sums have the bounds
\[
 \max_i\sum_jQ_{ij}^2+\norm Q_{\ell^4}^2
       +\norm{\diag Q}_2^2\leq C,
 \qquad \norm\Theta\leq C.
\]
Taylor expansion of $d_j\phi_i$ has linear parameter terms and a remainder
bounded entrywise by
\begin{equation}
 C\{Q_{ij}^2+Q_{ij}\Theta_j+\Theta_j^2\}.
 \label{eq:source-taylor-error}
\end{equation}
After multiplication by $s_i$, the linear parameter terms have bounded
nuclear norm: each is a rank-one matrix with a bounded-norm left vector
and a bounded-norm right vector.  For the $Q^2$ part, the matrix before
multiplication by $D_s$ has bounded Frobenius norm.  Its product with
$D_s$ has bounded nuclear norm by \eqref{eq:nuclear-tools} and
$\norm{D_s}_{\Frob}=\norm s\leq C$.
For the mixed part the sum of column norms is at most
\[
 C\sum_j\Theta_j\Bigl(\sum_i s_i^2Q_{ij}^2\Bigr)^{1/2}
 \leq C\norm\Theta
       \Bigl(\sum_i s_i^2\sum_jQ_{ij}^2\Bigr)^{1/2}\leq C.
\]
For the last part it is at most $C\norm s\sum_j\Theta_j^2\leq C$.
These separate estimates also apply to the combined remainder in
\eqref{eq:source-taylor-error}: entry by entry, assign each summand in
the bound its proportional share of the actual remainder.

The product rule also has a remainder with entries bounded by
\[
 C|d_js_i|\{Q_{ij}+\Theta_j\}.
\]
Assume $\norm{\nabla s}_{\Frob}\leq C$, which is zero for a seed,
holds for the initial implicit field by
Lemma~\ref{lem:implicit-construction}, and is inductive for later fields.
The $\Theta$ part has bounded sum of column norms by Cauchy--Schwarz.
For the $Q$ part, test against $P$ and split the trace into diagonal and
off-diagonal terms.  H\"older's inequality gives
\begin{align*}
 \sum_{i\ne j}|P_{ji}|\,|d_js_i|Q_{ij}
 &\leq \norm{P_{\rm off}}_{\ell^4}
       \norm{\nabla s}_{\Frob}\norm Q_{\ell^4}\leq C,\\
 \sum_i|P_{ii}|\,|d_is_i|Q_{ii}
 &\leq\norm P\norm{\diag\nabla s}_2\norm{\diag Q}_2\leq C.
\end{align*}
The first line uses exponents $4,2,4$. In particular, the argument does
not assume the diagonal $\ell^1$ bound for $\nabla s$ that we are still
proving. Bounded entries of $Q$ also give a bounded Frobenius norm for
this error. We have therefore controlled each source error both in
Frobenius norm and in trace against the indicated words.

\paragraph{Average corrections and size bounds.}
Next we control the averaged coefficients in the correction terms.
Differentiating a coefficient function in a primary argument gives another
source with the same inputs and bounds. Apply the size and Frobenius
parts of the preceding estimates to that source, then average:
\[
 |\av u|\leq\norm u/\sqrt n,\qquad
 \norm{d\av u}\leq\norm{\nabla u}_{\Frob}/\sqrt n,
\]
which proves the first line of \eqref{eq:small-average-bounds}.
The implicit coefficient $c_1$ already has these bounds from its direct
construction. For an auxiliary partial the scale is different: it is a
bounded coefficient vector, rather than a small vector. Averaging its
difference estimate
$|d_j\phi_i|\leq C(Q_{ij}+\Theta_j)$ and using the row bounds above
gives $|\av\phi|+\norm{d\av\phi}\leq C$.  This proves the second
line, including the self-coefficient $\chi$ when it occurs.

For a primary correction $\alpha_lm^{l-1}$ retain
$\alpha_l\nabla m^{l-1}$.  The dropped rank-one matrix has nuclear norm
at most $\norm{m^{l-1}}\norm{d\alpha_l}\leq C$.  The dropped product
error is $\nabla m^{l-1}$ times a diagonal of signed entries
$2d_i\alpha_l$.  Its nuclear norm is bounded by
\[
 2\norm{\nabla m^{l-1}}_{\Frob}\norm{d\alpha_l}
 \leq C\sqrt n\,n^{-1/2}\leq C.
\]
This also holds for $m^0=x$.  For an auxiliary correction $\gamma_bw^b$,
the analogous rank-one and product errors have bounded nuclear norm using
$\norm{w^b}+\norm{\nabla w^b}_{\Frob}\leq C$ and
$|\gamma_b|+\norm{d\gamma_b}\leq C$.
When primary derivatives in retained terms are replaced by
$H_l,M_l$, their Frobenius errors from
Lemma~\ref{lem:primary-derivatives} are multiplied either by a
small-vector diagonal or by an $O(n^{-1/2})$ scalar.  In the former case
\eqref{eq:nuclear-tools} applies; in the latter use
$\norm E_{\rm nuc}\leq\sqrt n\norm E_{\Frob}$.  Both give bounded
nuclear norm.

We can now close the induction for the size and Frobenius parts of
\eqref{eq:small-derivative-bound}. In \eqref{eq:small-field}, the terms are a
bounded-norm leading term, primary corrections bounded by
$C\sqrt n\,n^{-1/2}$, and bounded-norm auxiliary corrections.
The retained source derivative \eqref{eq:retained-source} is
Frobenius-bounded: use the preceding auxiliary derivative bounds and
$\norm{D_{s\partial_l\phi}\nabla h^l}_{\Frob}
\leq\norm{s\partial_l\phi}\norm{\nabla h^l}$.
The retained primary correction is bounded in Frobenius norm by
$|\alpha_l|\norm{\nabla m^{l-1}}_{\Frob}\leq C$.
Thus the size and Frobenius induction closes before we use or prove
the diagonal $\ell^1$ assertion.

\paragraph{Solving the differentiated self-dependence.}
The initial implicit source needs one additional step before the trace
cancellation. Retain its site partials, freeze the parameters and
averages, and replace the primary derivatives by their formal matrices.
This gives
\begin{equation}
 \begin{split}
 \nabla w&=A\nabla y^{\rm imp}+D_{\partial_k w}H_k+E_1,\\
 \nabla y^{\rm imp}&=(J-j\chi I)\nabla w-jc_1M_{k-1}+E_2.
 \end{split}
 \label{eq:implicit-formal-differentiation}
\end{equation}
The errors have bounded Frobenius norm and bounded trace against words
containing at most one $K$.  The parameter $c_1$ causes no exception:
write the corresponding source term as
$-j(\sqrt n c_1)A(t_*/\sqrt n)$, so that it is exactly of the source
form above.  Its difference bound follows from
\eqref{eq:implicit-frobenius}.
Solving \eqref{eq:implicit-formal-differentiation} yields
\begin{equation}
 \nabla w
 =K\{D_{\partial_k w}H_k-jc_1AM_{k-1}+E_1+AE_2\},
 \qquad K=(I-AJ+j\chi A)^{-1}.
 \label{eq:implicit-derivative-solved}
\end{equation}
Testing a propagated error against an ordinary word now puts at most
one $K$ into the test word. The preceding error estimates therefore give
both bounded trace and bounded Frobenius norm. Substitution gives the same
conclusion for
$\nabla y^{\rm imp}$. Subsequent fields are ordinary: propagating these
errors multiplies
them only by ordinary words. Induction therefore bounds the trace of
every final error against an ordinary word. Since the recipe is fixed
and finite, the required increases in word length are finite as well.

\paragraph{Cancellation of the formal traces.}
All differentiated errors are now controlled. To finish, we must bound
the diagonal $\ell^1$ norm of the retained formal derivative. A single
ordinary field already illustrates the cancellation. From
$w_i=\phi(h_i^1)s_i$ we obtain
$y=Jw-j\av{\phi'(h^1)s}x$. If $P=D_{\phi'(h^1)s}$, its retained
derivative is $JPJ-j\av{\phi'(h^1)s}I$, whose two diagonal predictions
cancel. We now show that the correction at each node gives this same
cancellation along every dependency path, including paths through the
implicit field.

At the point $x$, freeze the local partial vectors and all their averages. Index
each primary or auxiliary field by a node $a$, and attach its source to
the same node: for example, $h^l$ has source $m^{l-1}$, and $y^a$ has
source $w^a$. Let $p_{ad}$ be the partial of the source at node $a$
with respect to the field at node $d$. Introduce a formal variable $z$,
and denote the formal source and field derivatives by
$\mathfrak S_a$ and $\mathfrak H_a$, respectively. They obey
\begin{equation}
 \mathfrak S_a=\sum_dD_{p_{ad}}\mathfrak H_d,
 \qquad
 \mathfrak H_a=zJ\mathfrak S_a
             -z^2j\sum_d\av{p_{ad}}\mathfrak S_d.
 \label{eq:formal-graph}
\end{equation}
At the spin leaf, the source derivative for $h^1$ is $I$ and its field
derivative is $zJ$.  A terminal source uses just the first equation.
At $z=1$ these are precisely the retained formal derivatives.  In the
implicit case the only cyclic node is $y^{\rm imp}$ and its source $w$;
closing it gives
\begin{equation}
 \mathfrak S_{\rm imp}
 =K(z)\{D_{\partial_k w}\mathfrak H_{h^k}
                 -z^2jc_1A\mathfrak S_{h^k}\},
 \quad K(z)=(I-zAJ+z^2j\chi A)^{-1}.
 \label{eq:formal-self-loop}
\end{equation}
Every other dependency is acyclic and linear. Solving the single cyclic
node therefore leaves at most one inverse $K(z)$ in each product of the
closed formulas.

The averaged term in \eqref{eq:formal-self-loop} is essential even for
$k=2$. Write $D=D_{1-(m^1)^2}$ and $P=D_{\partial_2w}$, so that
$b_1=n^{-1}\Tr D$ and $c_1=n^{-1}\Tr P$. Then
\[
 \mathfrak H_{h^2}=z^2(JDJ-jb_1I),\qquad
 \mathfrak S_{h^2}=zDJ,
\]
and the implicit derivative is
$K(z)\{z^2P(JDJ-jb_1I)-z^3jc_1ADJ\}$.
At degree four, the two noncrossing pairings of $AJAJPJDJ$ have
predictions
\[
 j^2\chi b_1 AP,
 \qquad j^2c_1\frac{\Tr(AD)}n A.
\]
The first is canceled by the terms from $-j\chi A$ in the
quadratic coefficient of $K(z)$ and $-jb_1I$ in
$\mathfrak H_{h^2}$. The second is canceled by the degree-four
term $-jc_1AJADJ$ from the averaged correction. This calculation exhibits
the first cancellation involving both the
self-dependent node and its primary input. We now account for every
degree by following dependency paths.

We work coefficient by coefficient.  An uncontracted
dependency path, followed backwards from a field to the spin leaf, has
successive links $JD_{p_{ad}}$ and ends with the $J$ link of $h^1$.
A terminal source contributes an initial partial diagonal and then such
a path.  The choices of input nodes index these paths separately, preserving
all multiplicities.  Expanding \eqref{eq:formal-graph} is
inclusion--exclusion over disjoint originally adjacent link pairs: the
first term keeps the first link, while the second contracts it with the
next link, with factor $-j\av{p_{ad}}$, and continues from the input
node's source.  That continuation skips exactly the link just contracted.
Every contraction has degree two in $z$, equal to the degree of the
replaced links.  Consequently a coefficient of degree $L$ involves only
uncontracted paths of length $L$, of which there are finitely many even
with the self-dependence.

To compute the trace prediction, test the formal derivative by a bounded
diagonal $D$. Fix a nonempty path and a full noncrossing pairing of its
links. Reinserting contracted adjacent pairs restores that same
full-pairing prediction.
Each contraction contributes its minus sign.  Conversely, all subsets of
the originally adjacent pairs of that full pairing occur, with precisely
these signs.  Such a full pairing has at least one originally adjacent
pair, so its total coefficient is zero.  Paths of odd length have no full
pairing.  There is no empty path contributing a derivative of a small
seed, because all seeds are fixed with respect to $x$; parameter
derivatives were already put into the error. Thus the trace prediction
vanishes as a formal power series. The
polynomial-prediction assertion for single-inverse words in
Lemma~\ref{lem:words} then makes it zero at $z=1$ in the closed formulas.

To pass from this zero prediction to a bounded trace, we use the
small-vector structure. Every product in a closed formal derivative
of a small vector contains
either a small-vector diagonal or a scalar of size $O(n^{-1/2})$:
a path giving a nonzero derivative must enter the primary nodes before
reaching the spin leaf.  Entry from a small source uses
either its primary partial, a small vector, or the average of that partial
in a correction term.  The latter has the indicated scalar size.
For the implicit starting node these are exactly
$\partial_k w$ and $c_1$ in \eqref{eq:formal-self-loop}.

For a term containing a small diagonal $D_s$, rotate the tested trace
to put that factor first.  The word estimate for the remaining factors
and Cauchy--Schwarz give
\[
 \left|\Tr(D_sQ)-\sum_i s_i\,[\mathsf P(Q)]_{ii}\right|
 \leq\norm s\norm{\diag(Q-\mathsf P(Q))}_2\leq C.
\]
The normalized-trace prediction is cyclically invariant by
Lemma~\ref{lem:words}, so the rotation does not change the prediction
being cancelled.  For a term with a scalar of size $O(n^{-1/2})$, its
unscaled trace error is at most $C\sqrt n$ by the diagonal $\ell^2$
word estimate, and multiplication gives the same $O(1)$ bound.
The closed formula has finitely many terms. Combining their errors
with the zero prediction and the already controlled differentiated errors
gives
\[
 |\Tr(D\nabla U)|\leq C\qquad(\norm D\leq1,\ D\text{ diagonal}).
\]
Taking $D_{ii}$ to be the sign of $d_iU_i$ proves the remaining assertion
$\norm{\diag\nabla U}_1\leq C$.  All arguments were pointwise inside
the prescribed buffers, which proves the local version as well.
\end{proof}

\subsection{The weak residual rule}

We now apply the diagonal $\ell^1$ estimate to the exact conditional
integration-by-parts identity at the first primary field. The scalar
integral identity then reduces each higher residual to earlier ones.
We carry out this induction for ordinary tests and square-density tests
together.

\begin{lemma}[Weak residual rule]\label{lem:residual}
Let $U$ be a terminal vector in a fixed recipe admissible from level
$p\geq1$.  For an ordinary recipe and every scalar $G$,
\begin{equation}
 \left|\E G R_p^{\mathsf T}U\right|\leq C\norm G_*.
 \label{eq:weak-residual}
\end{equation}
For a local recipe with implicit starting index $k$, the same conclusion
holds when $p\leq k$ and $G$ is supported on $\Omega_\rho$.

There is also an ordinary square-density version.  If
$\sup_x(|H_0(x)|+\norm{dH_0(x)})\leq C_0$, then for every $f$,
\begin{equation}
 \left|\E f^2H_0 R_p^{\mathsf T}U\right|
       \leq C\{\E f^2+\D_{h_0}(f)\}.
 \label{eq:square-residual}
\end{equation}
Constants are uniform in $f,G,H_0$ subject to the stated bound, and may
depend on $p$, the recipe, its coefficient bounds, and the finite set of
word diagnostics needed in the proof.
\end{lemma}

\begin{proof}
We induct strongly on $p$ over all finite admissible recipes, proving
the two assertions together. Each step may enlarge the original recipe,
but it adds only finitely many auxiliary operations and invokes only
smaller residual indices. Consequently every particular instance of the
induction uses finitely many operations and a finite word-length bound.

At the base level $p=1$, condition on all spins except the one being
differentiated. Its conditional mean is $m_i^1$ and its conditional
variance is $v_i^1$. A function of this single spin is affine, with linear
coefficient $d_i$ of the function. Taking its conditional covariance
with the spin gives the exact identity
\begin{equation}
 \E\sum_i(x_i-m_i^1)GU_i
       =\E\sum_i v_i^1d_i(GU_i).
 \label{eq:conditional-ibp}
\end{equation}
By \eqref{eq:discrete-product}, its right side is
\[
 \E\sum_i v_i^1Gd_iU_i
   +\E\sum_i v_i^1(d_iG)\{U_i-2x_id_iU_i\}.
\]
The diagonal derivative bound makes the first term at most $C\E|G|$.
For the second term, the coefficient vector has pointwise norm at most
$\norm U+2\norm{\diag\nabla U}_2\leq C$.
Weighted Cauchy--Schwarz and $0<v_i^1\leq1$ therefore bound that term
by $C\D_{h_0}(G)^{1/2}$.  This proves \eqref{eq:weak-residual}.
For a supported local test, its derivatives vanish outside the one-flip
neighborhood of its support; all bounds used here are available there.

For \eqref{eq:square-residual}, absorb the bounded multiplier by replacing
$U$ in the base identity with $\wt U=H_0U$. This vector still has bounded
norm and bounded diagonal derivative
$\ell^1$ norm: the additional term obeys
$\sum_i|U_i d_iH_0|\leq\norm U\norm{dH_0}$, and the product-error
term is bounded by the same argument using
$\norm{\diag\nabla U}_2$.  Apply \eqref{eq:conditional-ibp} with
$G=f^2$ and use
\[
 d_i(f^2)=2f\,d_i f-2x_i(d_i f)^2.
\]
The diagonal-derivative term is bounded by $C\E f^2$; the remaining
terms are bounded by
$C\sqrt{\E f^2\,\D_{h_0}(f)}+C\D_{h_0}(f)$.
The base case is now proved in both forms. Before proceeding to higher
residuals, we record how bounded scalar multipliers affect the tests. If
$|\vartheta|+\norm{d\vartheta}\leq C$, then
\begin{equation}
 \norm{\vartheta G}_*\leq C'\norm G_*.
 \label{eq:star-multiplier}
\end{equation}
Indeed arrange the product rule as
$d_i(\vartheta G)=\vartheta(x^{(i)})d_iG+G(x)d_i\vartheta$,
square, and sum with weights $v_i^1\leq1$.  In the square-density
case such a multiplier can be absorbed into $H_0$; the product retains
its bounded size and difference norm.  For local tests only the buffer
bounds are needed, since multiplication does not enlarge their support.

For the induction step, suppose both assertions hold through level $p$
and take $U$ admissible from level $p+1$. The scalar integral identity
suggests the following new source. Set
\[
 a_p=j(b_p-b_{p-1}),\qquad
 U'_i=\Phi(h_i^p;h_i^{p+1},a_p)U_i.
\]
This is a source admissible from level $p$.  Its coefficient and all
needed derivatives are bounded because $|a_p|\leq j$ and the bounds for
$\Phi$ proved with \eqref{eq:phi-stein} hold uniformly in its two real
field arguments.
The level condition on $U$ makes its site partial at level $p$ vanish:
$\partial_pU$ differentiates only explicit primary arguments. Parameters
and auxiliary arguments remain fixed, including when their values depend
on earlier primary fields.  Applying \eqref{eq:phi-stein} coordinatewise
therefore yields
\[
 R_{p+1}^{\mathsf T}U
 =(h^p-h^{p+1}-a_pm^p)^{\mathsf T}U'
                      -a_p n\av{\partial_pU'}.
\]
The primary recursion gives
\[
 h^p-h^{p+1}-a_pm^p
       =JR_p+a_pR_p-jb_{p-1}R_{p-1},
\]
where the last term is absent for $p=1$.  Add a last auxiliary field
$y'$ by \eqref{eq:small-field} with source $U'$, and write
$\alpha_l=\av{\partial_lU'}$ and
$\gamma_b=\av{\partial_{y^b}U'}$.  Symmetry of $J$ now gives
\begin{align}
 R_{p+1}^{\mathsf T}U
 ={}&R_p^{\mathsf T}y'+a_pR_p^{\mathsf T}U'
              -jb_{p-1}R_{p-1}^{\mathsf T}U'
              +j\sum_b\gamma_bR_p^{\mathsf T}w^b\notag\\
    &+j\sum_{l\geq p}\alpha_lR_p^{\mathsf T}m^{l-1}
                   -a_pn\alpha_p.
 \label{eq:residual-induction-expansion}
\end{align}
The first line now consists of earlier residuals with admissible sources
and bounded multipliers. After testing, \eqref{eq:star-multiplier},
\eqref{eq:small-average-bounds}, and the level constraints allow us to
apply the induction hypothesis to each term. In particular the enlarged
recipe is admissible from level $p$, which is
also sufficient for its $R_{p-1}$ term.

On the second line, the terms with $l>p$ use the terminal source
$m^{l-1}/\sqrt n$, admissible from level $p$, and multiplier
$\sqrt n\alpha_l$. The multiplier bounds follow from
Lemma~\ref{lem:small-derivatives}. The remaining term is handled by the
identity
\begin{equation}
 R_p^{\mathsf T}m^{p-1}
       =n(b_p-b_{p-1})-R_p^{\mathsf T}m^p.
 \label{eq:residual-inner-product}
\end{equation}
This follows from
$(m^{p-1}-m^p)^{\mathsf T}(m^{p-1}+m^p)
=\norm{m^{p-1}}^2-\norm{m^p}^2$.
The first term on the right of \eqref{eq:residual-inner-product}
cancels $a_pn\alpha_p$ in
\eqref{eq:residual-induction-expansion}; the remaining term is handled
with the admissible source $m^p/\sqrt n$ and multiplier
$\sqrt n\alpha_p$. This completes the induction for both
\eqref{eq:weak-residual} and
\eqref{eq:square-residual}. In the local version it stops at level $k$;
each added source respects the support and buffer requirements of the
initial implicit construction.
\end{proof}

The last step puts the implicit source into the form needed for the
posterior comparison. The following identity is a direct application of
the weak residual rule, with the averaged correction retained.

\begin{lemma}[Terminal residual identity]\label{lem:terminal-residual}
For the implicit construction at index $k$ and every $G$ supported on
$\Omega_\rho$,
\begin{equation}
 \left|\E G\left[
 \{h^k-h_0-(J-jb_{k-1}I)m^k\}^{\mathsf T}w
                   -j(b_k-b_{k-1})n c_1\right]\right|
       \leq C\norm G_*.
 \label{eq:terminal-residual}
\end{equation}
\end{lemma}

\begin{proof}
The vector in braces is $JR_k-jb_{k-1}(R_{k-1}+R_k)$.
To express its pairing with the implicit source in terms of earlier
residuals, use symmetry of $J$, the equation
$Jw=y^{\rm imp}+j\chi w+jc_1m^{k-1}$, and
\eqref{eq:residual-inner-product}. The bracketed expression becomes
\[
 R_k^{\mathsf T}y^{\rm imp}
   +j(\chi-b_{k-1})R_k^{\mathsf T}w
   -jb_{k-1}R_{k-1}^{\mathsf T}w
   -jc_1R_k^{\mathsf T}m^k.
\]
Apply Lemma~\ref{lem:residual} to the first three terms, absorbing their
bounded parameters into the test. For the last term, use source
$m^k/\sqrt n$ and multiplier $\sqrt n c_1$. The bounds needed for this
multiplier are precisely \eqref{eq:implicit-size} and
\eqref{eq:implicit-frobenius}.
\end{proof}
